%! Characteristics of Functions: Increasing, Decreasing, Extrema, and Concavity — Unit 2, Lesson 2.5 — Warm-Up
% Exactly ONE page, blank and keyed (LESSON_SHAPE.md, one_page). Four prerequisite items:
% 1 is interval notation (every answer today is an interval); 2 reads a graph in BOTH
% directions — an output from an input, then inputs from an output — and re-states 2.4's
% domain and range; 3–4 are slope and its sign, the line-sized version of rising and falling.
\documentclass[10pt]{article}
\usepackage{precalculus-article}
\usepackage{precalculus-boxes}
\newcommand{\TallMath}[1]{$\displaystyle #1\rule[-1.4em]{0pt}{3.2em}$}

\begin{document}

\pageheader{Unit 2, Lesson 2.5}{Warm-Up}

\vspace{0.12in}

\begin{enumerate}[itemsep=12pt]

\item Write each set of numbers in interval notation.

\begin{enumerate}[label=(\alph*), itemsep=8pt]
  \item All $x$ with $2 < x < 5$. \hfill \blank{3cm}
  \item All $x$ with $-1 \le x \le 4$. \hfill \blank{3cm}
\end{enumerate}

\item \begin{minipage}[t]{0.52\linewidth}
The graph of $g$ is shown, including both endpoints.

\begin{enumerate}[label=(\alph*), itemsep=10pt]
  \item Find $g(0)$. \hfill \blank{2.2cm}
  \item For which $x$ is $g(x) = -2$? \hfill \blank{2.2cm}
  \item State the domain of $g$ in interval notation. \hfill \blank{2.2cm}
  \item State the range of $g$ in interval notation. \hfill \blank{2.2cm}
\end{enumerate}
\end{minipage}\hfill
\begin{minipage}[t]{0.4\linewidth}
\vspace{-6pt}
\centering
\begin{tikzpicture}[x=0.5cm, y=0.5cm, baseline=(current bounding box.north)]
  \draw[linegray, very thin, step=1] (-3,-3) grid (6,5);
  \draw[->, thick] (-3,0) -- (6.5,0) node[right] {\small $x$};
  \draw[->, thick] (0,-3) -- (0,5.5) node[above] {\small $g(x)$};
  \foreach \x in {-2,-1,1,2,3,4,5} \node[below, font=\scriptsize] at (\x,0) {\x};
  \foreach \y in {-2,-1,1,2,3,4} \node[left, font=\scriptsize] at (0,\y) {\y};
  \draw[very thick, plum] (-2,1) -- (0,4) -- (3,-2) -- (5,2);
  \fill[plum] (-2,1) circle (2.5pt);
  \fill[plum] (5,2) circle (2.5pt);
\end{tikzpicture}
\end{minipage}

\item Find the slope of the line through $(1, 3)$ and $(4, 9)$.

\begin{work}
  m &= \frac{9 - 3}{4 - 1} \\
    &= \frac{6}{3} = 2
\end{work}

Slope: \blank{2.2cm}

\item A line falls from left to right. Is its slope positive or negative? As $x$ gets larger,
what happens to $y$?

\vspace{4pt}
\blank{9cm}

\end{enumerate}

\end{document}
